% ============================================================
% v2 — Capitolo 3: Il progetto — metodo e campagna
% Blocco II di IV. Budget: ~4 pagine.
% Include, condensati: regole di misura, protocollo a gate,
% i tre episodi di validità dell'ambiente (con il pilota di aprile
% ridotto a un paragrafo) e il design della campagna 2x2.
% ============================================================
\chapter{Il progetto --- metodo e campagna}
\label{ch:metodo}

La piattaforma del \cref{ch:piattaforma} non è stata progettata in un
solo passaggio: è il prodotto di una campagna iterativa in cui ogni
modifica doveva superare un gate quantitativo predefinito prima di
entrare nel trunk. Questo capitolo documenta le regole di misura
(\cref{sec:metriche}), il protocollo di gate con il registro delle
iterazioni (\cref{sec:gate}), i difetti di validità trovati e corretti
lungo il percorso (\cref{sec:validita}) e il design della campagna di
confronto finale (\cref{sec:campagna}).

% ------------------------------------------------------------
\section{Metriche e regole di misura}
\label{sec:metriche}

\paragraph{Unità di misura e integrità.}
L'unità di analisi è il \emph{record a livello di agente}: una voce JSON
per agente per episodio, sei per episodio. Tutti gli aggregati sono
calcolati dopo deduplica per identificatore di episodio e di agente ---
in valutazione lo scenario fa parte della chiave. Due condizioni di
integrità sono imposte prima che qualunque numero venga riportato:
esattamente sei record per episodio, e nessun valore non finito. Una run
che fallisce una delle due non viene analizzata oltre.

\paragraph{Criterio di successo rigoroso.}
Un agente ha successo se e solo se la sua ragione di terminazione è
\texttt{route\_complete}; in particolare \texttt{route\_short} ---
raggiungere la fine di un percorso degenere troppo corto --- \emph{non}
è pass. La regola è stata messa alla prova durante lo sviluppo, dove
una modifica che contava \texttt{route\_short} come successo è stata
introdotta e poi revertita, mantenendo il criterio rigoroso al costo di
numeri assoluti più bassi. Questo approccio è risultato un pregio architetturale
poiché le metriche sono risultate più accurate e robuste.

\paragraph{Tre viste di reporting.}
Ogni risultato è riportato per veicoli e pedoni separatamente, più una
vista combinata: le due classi hanno task, spazi di azione e limiti
strutturali diversi, e il registro mostra effetti di regime che esistono
per una sola classe. Il successo congiunto a livello di episodio non è
mai usato, perché con agenti correlati compone i fallimenti in modo
moltiplicativo e nasconde la struttura di classe. I valori cumulativi
sono l'aggregazione primaria, il successo a finestra scorrevole è usato
solo dal gate del curriculum, e la dinamica entro la run è mostrata
dividendo l'addestramento in quartili (Q1--Q4). La \cref{tab:metriche}
definisce le metriche usate nel \cref{ch:risultati}.


\begin{table}[t]
  \centering\small
  \caption{Catalogo delle metriche. Tutti i tassi sono quote dei record
  a livello di agente.}
  \label{tab:metriche}
  \begin{tabular}{@{}p{3.4cm}p{9.4cm}@{}}
    \toprule
    Metrica & Definizione \\
    \midrule
    Tasso di successo (SR) & record con \texttt{route\_complete} \\
    Collisione / offroad & record con \texttt{collision} / \texttt{offroad} (offroad solo veicoli) \\
    Stuck / timeout & record troncati con completamento $< 0.3$ o loop attivo / con completamento $\geq 0.3$; la somma s$+$t è il target primario anti-stallo \\
    Completamento / efficienza di percorso & frazione di waypoint consumati; lunghezza ottima completata su distanza percorsa, con tetto a 1 e valore 0 per i record stuck \\
    Velocità media & km/h per i veicoli, m/s per i pedoni \\
    Diagnostica dei percorsi & quote di sorgente, flag under-target e too-short, attraversamenti, passi di non-progresso \\
    \bottomrule
  \end{tabular}
\end{table}


% ------------------------------------------------------------
\section{Protocollo di sviluppo guidato da gate}
\label{sec:gate}

La varianza da run a run in questo ambiente è risultata elevata e il budget di
calcolo non permetteva di replicare ogni modifica candidata. Il
protocollo adottato in risposta ha quattro regole. \textbf{(1) Finetuning sequenziale}: 
ogni candidato è isolato comportamentalmente --- un singolo termine di reward, 
una singola soglia, un singolo gruppo di osservazioni. 
\textbf{(2) A/B appaiato a budget diagnostico fisso}: 
ogni candidato gira per $3\times10^{5}$ passi contro la baseline 
corrente del trunk, con lo stesso seed e --- grazie al
seeding deterministico dei percorsi (\cref{sec:mondo}) --- una sequenza
di percorsi appaiata per costruzione; le run di sviluppo erano bloccate
su easy salvo quando il candidato riguardava il curriculum stesso.
\textbf{(3) Un gate
quantitativo predefinito}: per le modifiche orientate ai veicoli il
candidato è promosso solo se, rispetto alla baseline appaiata, l'SR
veicoli migliora di almeno $+2.0$\,pp, stuck$+$timeout migliora di
almeno $-2.0$\,pp, collisione e offroad peggiorano ciascuno di non più
di $+1.0$\,pp, ed entrambe le condizioni di integrità valgono; le
modifiche che toccano macchinari condivisi portano una condizione
laterale \emph{must-not-collapse} sull'altra classe, e un gate fallito
significa revert. \textbf{(4) Meccanismo ed esito registrati
separatamente}: il registro annota se il meccanismo inteso si è
verificato e se il gate comportamentale è passato --- i due spesso
dissentono, e ed alcuni finetuning che hanno fallito i gate sono stati 
mantenuti consapevolmente con l'ipotesi registrata come falsificata,
perché revertirle avrebbe ripristinato un difetto o perché il
fallimento stesso era informativo.

Lo sviluppo è proceduto in sei fasi; il registro completo dei candidati
è nell'\cref{app:registro}. Un \emph{pilota architetturale} di nove run
(aprile 2026) ha incrociato encoder del critic e regimi nell'era
pre-correzioni: i dati erano integri, ma i veicoli restavano a un tasso
di successo prossimo a zero (0.2--3.9\%) in ogni cella, rendendo il
confronto non informativo e reindirizzando lo sforzo verso la validità
dell'ambiente. Sono seguite una fase di \emph{correttezza}, una di
\emph{osservabilità} (feature veicolo consapevoli del percorso, poi
l'estensione a 47D che chiude il gap di aliasing di stato), una
sequenza di \emph{iterazioni su ottimizzatore e reward} dagli esiti
misti, una fase sull'\emph{accoppiamento tra classi} --- pedoni più
veloci saturano il canale di hazard dei veicoli e possono spingerli in
un equilibrio difensivo a velocità zero attraverso il gate
\texttt{safe\_to\_push}, risolto alzandone la soglia da 0.75 a 0.85 ---
e infine la \emph{meccanica del curriculum}, con l'ultima modifica
entrata nel trunk prima della campagna: il tetto agli attraversamenti
pedonali portato da otto a tre.

Un ultimo dato di questa fase è il metro di tutta la tesi. Tre repliche
a lunghezza piena ($3\times10^{6}$ passi) dell'\emph{identico} trunk
pre-campagna, con codice, configurazione e seed byte-identici, hanno
prodotto tassi di successo cumulativi dei veicoli tra 52.2 e 62.4\%:
CARLA sotto carico multi-agente non è deterministico da run a run
nemmeno in modalità sincrona con tutti i seed fissati. Questa banda di
$\sim$10\,pp è il riferimento contro cui ogni effetto del
\cref{ch:risultati} viene letto, e motiva sia il design appaiato sia la
cautela attribuita ai delta a replica singola.

% ------------------------------------------------------------
\section{Validità dell'ambiente: bug trovati e corretti}
\label{sec:validita}

In una piattaforma multi-agente custom l'apparato di misura è parte
dell'esperimento. Tre episodi durante lo sviluppo hanno prodotto
apparenti effetti di policy che erano in realtà artefatti; sono
riportati qui sia perché hanno plasmato il protocollo finale, sia perché
le loro firme di rilevamento possono essere utili ad altri.

\paragraph{Lunghezza dei percorsi: un cambio di distribuzione dei task.}
Il route planner poteva emettere percorsi molto più lunghi del target di
livello, quindi episodi nominalmente ``easy'' spesso non erano facili.
Imporre i limiti $[0.5, 2.0]\times$ il target (\cref{sec:percorsi}) ha
spostato l'SR cumulativo di addestramento dei veicoli di
$\approx +50$\,pp \emph{senza alcuna modifica alla policy o alla
reward}, e il pavimento dei veicoli del pilota di aprile è stato poi
ricondotto proprio a questo difetto. È l'evidenza interna più forte per
una regola usata in tutta la tesi: \emph{mai confrontare numeri assoluti
tra versioni dell'ambiente; ancorare le affermazioni a contrasti
appaiati entro una versione}.

\paragraph{Seeding dei percorsi: appaiamento rotto.}
I seed dei percorsi per agente derivavano originariamente dalla
\texttt{hash()} built-in di Python: due run nominalmente identiche 
estraevano sequenze di percorsi \emph{diverse}, degradando silenziosamente 
ogni confronto A/B. La correzione deriva i seed da una \texttt{SeedSequence} 
esplicita (\cref{sec:mondo}); la lezione è che l'appaiamento va costruito e 
verificato, non assunto.

\paragraph{Pipeline di valutazione: collasso della misura.}
La prima valutazione di checkpoint ha restituito un tasso di successo
esattamente del 33.33\% in \emph{ogni} scenario: uno dei tre veicoli a
completare, ogni episodio. Questo pattern troppo pulito per essere vero
ha esposto quattro difetti accoppiati: il seed di scenario per episodio
era costante, quindi ogni episodio rigiocava la stessa configurazione di
traffico; la policy era valutata sulla sua media deterministica,
collassando la gaussiana su una singola traiettoria; il sottoprocesso
non inizializzava mai il logger; e il seed torch era costante tra i
sottoprocessi. Tutti e quattro sono stati corretti e il protocollo
stocastico della \cref{sec:critic} è diventato lo standard; nessuna di
queste correzioni tocca l'addestramento, e nessuna affermazione di
policy è derivata da esse. La lezione generale: la validità della misura
va stabilita \emph{prima} del tuning, e risultati implausibilmente
regolari sono un segnale di rilevamento, non una buona notizia.

% ------------------------------------------------------------
\section{Design della campagna finale}
\label{sec:campagna}

L'esperimento finale è un fattoriale $2 \times 2$ --- \{encoder del
critic MLP, GNN\} $\times$ \{curriculum, mixed-batch\} --- di quattro
run da $3\times10^{6}$ passi ciascuna, tutte sullo stesso trunk
congelato, con un unico seed (999) appaiato tra le quattro celle e la
difficoltà definita dalla lunghezza del percorso. Le configurazioni sono
byte-identiche tra le celle salvo i due flag sperimentali. Il design era
originariamente dimensionato a tre seed appaiati per cella; l'asse di
replicazione è stato eliminato per ragioni di budget di calcolo, poiché
ogni run occupa l'unica macchina disponibile per circa due giorni
(\cref{app:riproducibilita}). La riduzione è dichiarata qui anziché
nascosta: con un seed per cella, ogni contrasto del \cref{ch:risultati}
è un \emph{caso di studio appaiato}, non la stima di una distribuzione.

\paragraph{Contrasti.}
Il contrasto primario è curriculum contro mixed-batch \emph{entro}
architettura, ed è osservato \emph{due volte} --- una per ciascuna
architettura del critic, addestrata indipendentemente --- quindi la
replicazione tra architetture sostituisce in parte quella sui seed: un
effetto che compare con lo stesso segno in entrambe le coppie è molto
più difficile da attribuire al rumore di quanto lo sarebbe un contrasto
singolo. Il contrasto secondario è critic MLP contro GNN \emph{entro}
regime; la lettura dell'interazione (H3) è solo direzionale a questa
dimensione campionaria. Nessuna inferenza statistica formale è
rivendicata: gli effetti sono letti contro la banda delle repliche a
seed identico della \cref{sec:gate}.

La campagna \emph{non} mette invece alla prova: famiglie di difficoltà
basate sulla densità di traffico o su assi misti; l'encoder del critic ad
attenzione e la normalizzazione PopArt (disponibili, disabilitati);
architetture degli actor (sempre MLP); ablazioni del numero di agenti
(sempre $3+3$); o coppie di mappe oltre Town03 $\to$ Town05.
